\begin{papercard}{macqueen1998}
\cf{Stages} \stagetag{C} \stagetag{I}. The standard practitioner paper on \emph{codebook structure and team-based refinement}: a six-part codebook entry, a refinement loop built on independent coding of the same sample, and eight practical lessons. Its source is CDC team-based qualitative projects (interviews and open-ended survey answers).

\cf{Why a structured codebook} Codes are ``the building blocks for theory or model building'' and embody the analyst's assumptions \pc{macqueen1998}{1}{Codes are the building blocks for theory or model building}. CDC uses multiple coders to assess reliability and validity through intercoder agreement and to handle volume \pc{macqueen1998}{1}{We generally use multiple coders so that we can assess the reliability and validity of the coded data}. The paper covers (1) the codebook as a stable frame, (2) features that improve agreement, (3) the value of team development \pc{macqueen1998}{1}{This paper describes (1) how a structured codebook provides a stable frame}.

\cf{Origin of the format} It grew out of coding processes-of-change in about 600 HIV-prevention interviews. A coding manual by O'Connell added a section ``differentiating'' a code from others. This was reframed as \emph{exclusion criteria} alongside \emph{inclusion criteria} \pc{macqueen1998}{1}{They conceptu- alized the "differentiation" component as "exclusion criteria," which complemented}.

\cf{Six-component codebook entry} Code (mnemonic), brief definition, full definition, when to use, when \emph{not} to use, examples \pc{macqueen1998}{2}{The codebook structure has evolved to include six basic components: the code, a brief definition}. Table 1 shows a worked entry (MARGIN: marginalised community members) whose ``When Not to Use'' points to related codes INSTIT and SELFMAR. A relational database makes the structure a template, allows brief and full versions, change reports by date field, export to analysis software, and hierarchical code families \pc{macqueen1998}{2}{There are practical advantages to using a database management program for codebook development}.

\cf{The codebook as a frame} The codebook ``forces the analyst to place his or her assumptions and biases in plain view''. The adequacy of answers depends on the sensitivity and specificity of the codes and the reliability of links to text \pc{macqueen1998}{3}{It forces the analyst to place his or her assumptions and biases in plain view}.

\cf{Coding process} Index codes tag text for retrieval and measurement codes assign values (presence/absence, amount). Either way, explicit rules for segment boundaries are needed \pc{macqueen1998}{3}{Thus, the coding process must include explicit quidelines for defining the boundaries of the text}. For long interviews: Stage 1 \emph{structural coding} by research question, keeping the full elicitation and response in each segment. For in-depth coding, prefer \emph{emic} (respondents' own terms) over \emph{etic} jargon codes, since the latter smuggle in coder preconceptions. Example: three class codes were replaced by one income/class code \pc{macqueen1998}{3}{We have found that code definitions that incorporate substantial references to professional jargon}. For short survey answers: a 0/1 presence matrix by participant $\times$ code (Table 2), with hierarchical codes so categories can be collapsed later instead of limiting the codebook up front \pc{macqueen1998}{4}{Such a strategy is preferable to limiting the content ofthe codebook a priori}.

\cf{Refining the codebook (Figure 1)} (1) An initial code list from etic concepts, emic themes or both, led by one or two members and reviewed by the team. (2) Definitions proposed and reviewed. (3) Two or more coders \emph{independently code the same sample}. (4) Compare segmentation and code application; if acceptable, continue with periodic checks, and if not, review. Inconsistencies are triaged: \emph{coder error} (a training issue) vs.\ \emph{code-definition problems} such as overlapping inclusion criteria (a team issue). After clarifying the codebook, all previously coded text is reviewed and recoded \pc{macqueen1998}{5}{Once the problems are identified and the codebook clarified, all previously coded text is}. Agreement is assessed by segment-by-segment review (boundaries and codes), and kappa. Kappa spots coder error unless all coders are mistrained, and flags poorly defined codes that are applied erratically \pc{macqueen1998}{5}{It is also helpful for identifying poorly defined codes, which tend to be applied erratically}. A caution: subsample comparison tends to miss text no coder captured \pc{macqueen1998}{5}{In particular, it is prone toward the loss of pertinent text}.

\cf{Eight practical lessons} (1) One person owns the codebook and clears all coding questions. (2) Review every code regularly, since implicit rules drift. (3) A codebook someone can learn in a few hours, with coders handling 30--40 codes at a time, staged by domain beyond 40 \pc{macqueen1998}{5}{For the most part, coders can reasonably handle 3 0 4 0 codes at one time}. (4) Write the segmentation plan while writing the codebook. (5) Fix agreement measures \emph{early} and decide a priori the permitted variability. (6) Assume nothing is obvious and define abbreviations. (7) Delete deadwood codes, using one generic UNIQUE code for single-response codes. (8) Accept recoding as ``forward movement'' \pc{macqueen1998}{6}{Finally, accept the fact that text will need to be recoded as the codebook is refined}.

\cf{Use for this thesis} The six-field entry is a ready template for the thesis codebook (add \emph{observable evidence} and \emph{screenshot reference} fields). The exclusion criteria and cross-references to neighbouring codes target exactly the confusion between similar constructs (e.g.\ consent vs.\ notice). The two-track triage of disagreements (coder error vs.\ definition problem) gives a defensible rule for when to revise the codebook vs.\ retrain. Plan agreement thresholds \emph{before} piloting. Because it is a methods note, not an empirical validation, cite it for the structure and rationale, not as proof of reliability.
\end{papercard}
